\cotefile{136}
\scanpage{95}
\begin{center}- \mspage{169} -\end{center}
\margin{IX}

finies ; donc $X\longmapsto\widehat X$ transforme groupes en groupes. Notons que $\widehat X$ ne dépend en fait que de la restriction de $X$ à un voisinage ouvert arbitraire $U$ de $S$.

\uline{Lemme} 13.2.1. \uline{Les données sont celles de} 13.1, \uline{et on suppose que les conditions suivantes sont vérifiées} : a) $A=A^0$, b) $R$ \uline{est constant au voisinage de} $s$, \uline{et} c) $R_s$ \uline{engendre l'espace vectoriel complexe} $E_s$. \uline{Soit} $A'$ \uline{un deuxième groupe analytique lisse sur} $S$, \uline{et considérons l'application}
\[
u\longmapsto\widehat u:\operatorname{Hom}_{S\text{-}\mathrm{gr}}(A,A')\longrightarrow\operatorname{Hom}_{\widehat S\text{-}\mathrm{gr}}(\widehat A,\widehat {A'})\quad,
\]
\uline{où les notations sont celles de} 13.2, \uline{et l'application analogue obtenue en se restreignant au-dessus d'un voisinage ouvert} $U$ \uline{variable de} $s$, \uline{et passant à la limite inductive} :
\[
(13.2.1.1)\qquad u\longmapsto\widehat u:\operatorname{Hom}_{(S,s)\text{-}\mathrm{gr}}(A,A')\longrightarrow\operatorname{Hom}_{\widehat S\text{-}\mathrm{gr}}(\widehat A,\widehat {A'})\quad,
\]
(\uline{où} $(S,s)$ \uline{désigne le} germe \uline{d'espace analytique de} $S$ \uline{en} $s$). \uline{Alors l'application précédente est} bijective.

Considérons en effet la suite exacte
\diagram{\begin{tikzcd}
0\arrow[r]&R'\arrow[r]&E'\arrow[r]&A'{}^0\arrow[r]&0
\end{tikzcd}}
analogue à (13.1.3) définie par $A'$. Compte tenu du fait que les homomorphismes $A\longrightarrow A'$ se factorisent nécessairement par $A'{}^0$, et de même après les changements de base $U\longrightarrow S$ et $S_n\longrightarrow S$, on peut supposer $A'=A'{}^0$. Compte tenu du dictionnaire 13.1.4, il faut prouver que la donnée d'un

\margin{\unsure{487}\ \illegible{2}}
\margin{487}
\margin{94}
